% ==============================================================================
% styles/layout.tex
% ==============================================================================

% ------------------------------------------------------------------------------
% 1. هندسه صفحه (Geometry)
% ------------------------------------------------------------------------------
% اعمال تنظیمات حاشیه‌ها بر اساس انتخاب کاربر در فایل settings.tex
\ifthenelse{\equal{\TheMarginType}{normal}}{
    \geometry{a4paper, top=3.5cm, bottom=3.5cm, inner=2.5cm, outer=3.5cm, headheight=40pt, headsep=1.2cm, footskip=1.5cm}
}{}
\ifthenelse{\equal{\TheMarginType}{narrow}}{
    \geometry{a4paper, top=3.5cm, bottom=2.5cm, inner=1.5cm, outer=1.5cm, headheight=40pt, headsep=1.2cm, footskip=1cm}
}{}
\ifthenelse{\equal{\TheMarginType}{wide}}{
    \geometry{a4paper, top=4cm, bottom=4cm, inner=3cm, outer=5cm, headheight=40pt, headsep=1.5cm, footskip=1.5cm}
}{}

% ------------------------------------------------------------------------------
% 2. سربرگ و پابرگ (Fancyhdr)
% ------------------------------------------------------------------------------
\usepackage{fancyhdr}
\pagestyle{fancy}
\fancyhf{} % پاک کردن پیش‌فرض‌ها

% استایل هدر (سربرگ)
\fancyhead[RO, LE]{\bfseries\textcolor{ThemePrimary}{\thepage}} % شماره صفحه
\fancyhead[LO]{\small\textcolor{ThemePrimary}{\rightmark}} % نام بخش در سمت درون (فرد)
\fancyhead[RE]{\small\textcolor{ThemePrimary}{\leftmark}}  % نام فصل در سمت درون (زوج)

% استایل فوتر (پابرگ)
% پابرگ خالی است تا فضای کافی برای زیرنویس‌ها (Footnotes) فراهم شود.
\fancyfoot{}

% خط کشی ملایم برای سربرگ
\renewcommand{\headrulewidth}{1pt}
\renewcommand{\headrule}{\hbox to\headwidth{\color{ThemePrimary}\leaders\hrule height \headrulewidth\hfill}}
\renewcommand{\footrulewidth}{0pt}

% استایل صفحات شروع فصل و فهرست
\fancypagestyle{plain}{
    \fancyhf{}
    \fancyhead[RO, LE]{\bfseries\textcolor{ThemePrimary}{\thepage}} % شماره صفحه
    \fancyhead[LO]{\small\textcolor{ThemePrimary}{\rightmark}} % نام بخش در سمت درون (فرد)
    \fancyhead[RE]{\small\textcolor{ThemePrimary}{\leftmark}}  % نام فصل در سمت درون (زوج)
    \renewcommand{\headrulewidth}{1pt}
    \renewcommand{\headrule}{\hbox to\headwidth{\color{ThemePrimary}\leaders\hrule height \headrulewidth\hfill}}
}

% ------------------------------------------------------------------------------
% 3. استایل‌دهی به تیترها (Titlesec)
% ------------------------------------------------------------------------------
% سفارشی‌سازی بخش‌ها با خط رنگی
\titleformat{\section}[block]
  {\needspace{5\baselineskip}\LARGE\bfseries\color{ThemePrimary}}
  {\thesection}
  {1em}
  {#1}
  [{\vspace{2pt}\color{ThemeSecondary!50}\titlerule[1.5pt]}]

\titleformat{\subsection}
  {\needspace{3\baselineskip}\Large\bfseries\color{ThemePrimary}}
  {\thesubsection}
  {1em}
  {#1}

\titleformat{\subsubsection}
  {\needspace{2\baselineskip}\large\bfseries\color{ThemeSecondary}}
  {\thesubsubsection}
  {1em}
  {#1}

% ------------------------------------------------------------------------------
% 4. تنظیمات فهرست مطالب (TOC)
% ------------------------------------------------------------------------------
% استفاده از tocloft برای زیبایی بیشتر فهرست‌ها
\usepackage{tocloft}

% وسط‌چین کردن و تغییر استایل عنوان فهرست
\renewcommand{\cfttoctitlefont}{\hfill\Huge\bfseries\color{ThemePrimary}}
\renewcommand{\cftaftertoctitle}{\hfill}

% استایل خطوط راهنما (نقطه‌چین) ملایم‌تر
\renewcommand{\cftdot}{\ensuremath{\cdot}}
\renewcommand{\cftsecdotsep}{4.5}
\renewcommand{\cftchapdotsep}{\cftdotsep}

% فاصله‌ها در فهرست
\setlength{\cftbeforechapskip}{1.5em}
\setlength{\cftbeforesecskip}{0.3em}

% تنظیم رنگ‌ها و فونت‌های فهرست مطالب هماهنگ با تیترها
\renewcommand{\cftchapfont}{\bfseries\large\color{black}}
\renewcommand{\cftchappagefont}{\bfseries\large\color{black}}

\renewcommand{\cftsecfont}{\color{ThemePrimary}}
\renewcommand{\cftsecpagefont}{\color{ThemePrimary}}

\renewcommand{\cftsubsecfont}{\color{ThemeSecondary}}
\renewcommand{\cftsubsecpagefont}{\color{ThemeSecondary}}

% تو رفتگی شماره صفحات در فهرست از سمت چپ (RTL)
\renewcommand{\cftchapafterpnum}{\hspace*{1.5em}}
\renewcommand{\cftsecafterpnum}{\hspace*{1.5em}}
\renewcommand{\cftsubsecafterpnum}{\hspace*{1.5em}}

% ------------------------------------------------------------------------------
% 4. تنظیمات فاصله خطوط و پاراگراف‌ها
% ------------------------------------------------------------------------------
% این تنظیمات به خوانایی متن فارسی بسیار کمک می‌کنند
\linespread{1.35} % فاصله خطوط بهینه برای متون فارسی
\setlength{\parskip}{0.6\baselineskip plus 2pt minus 1pt} % فاصله بین پاراگراف‌ها
\setlength{\parindent}{1.5em} % تورفتگی خط اول پاراگراف

% ------------------------------------------------------------------------------
% 5. ارجاع‌دهی پیشرفته (Hyperref & Cleveref)
% ------------------------------------------------------------------------------
% پکیج hyperref برای لینک‌های داخلی (رنگ‌های ملایم و حرفه‌ای)
\usepackage[
    colorlinks=true,
    linktoc=all,
    linkcolor=ThemePrimary, % رنگ آبی تیره و ملایم برای لینک‌های داخلی
    urlcolor=ThemePrimary,  % لینک‌های وب
    citecolor=ThemeSecondary % رنگ برای ارجاعات مراجع
]{hyperref}

% ساخت بوک‌مارک‌های تمیز در فایل PDF خروجی
\usepackage{bookmark}

% پکیج cleveref برای ارجاع‌دهی هوشمند (باید حتماً بعد از hyperref لود شود)
\usepackage{cleveref}

% سفارشی‌سازی نام‌های فارسی برای ارجاع هوشمند
\crefname{equation}{معادله}{معادلات}
\crefname{figure}{شکل}{اشکال}
\crefname{table}{جدول}{جداول}
\crefname{chapter}{فصل}{فصول}
\crefname{section}{بخش}{بخش‌ها}
\crefname{theorem}{قضیه}{قضایا}
\crefname{lemma}{لم}{لم‌ها}
\crefname{corollary}{نتیجه}{نتایج}
\crefname{definition}{تعریف}{تعاریف}
\crefname{example}{مثال}{مثال‌ها}
